%% =============================================================================
%% APIP -- full (long) version, flattened into a single self-contained file.
%% Generated from main.tex + sections/*.tex + main.bbl; no \input, no BibTeX.
%% Build with: pdflatex APIP_full.tex  (run twice for cross-references)
%% Source of truth remains main.tex + sections/ in the research_papers repo.
%% =============================================================================

% =============================================================================
% Your Agent is Underperforming: The Case for Agent Performance Improvement Plans
%
% Source: research/apip/APIP_Paper_v8_tracked.docx (tracked changes accepted)
% Style:  IEEE conference template (ieee_agc/latex) + IEEEtran BibTeX
%         style (ieee_agc/bibtex)
%
% Compilation:
%   pdflatex main.tex
%   bibtex   main
%   pdflatex main.tex
%   pdflatex main.tex
%
% Structure:
%   main.tex                     -- root document (this file)
%   sections/abstract.tex        -- abstract + index terms
%   sections/introduction.tex    -- Sec. 1
%   sections/background.tex      -- Sec. 2
%   sections/taxonomy.tex        -- Sec. 3
%   sections/framework.tex       -- Sec. 4
%   sections/schema.tex          -- Sec. 5
%   sections/simulation.tex      -- Sec. 6
%   sections/open_problems.tex   -- Sec. 7
%   sections/conclusion.tex      -- Sec. 8
%   references.bib               -- bibliography entries
% =============================================================================

\documentclass[conference]{IEEEtran}
\IEEEoverridecommandlockouts

% ---------------- Packages ----------------
\usepackage[T1]{fontenc}
\usepackage[utf8]{inputenc}
\usepackage{cite}
\usepackage{amsmath,amssymb,amsfonts}
\usepackage{graphicx}
\usepackage{textcomp}
\usepackage{xcolor}
\usepackage{booktabs}
\usepackage{array}
\usepackage{url}

% Fixed-width paragraph cells for the taxonomy/schema/simulation tables:
% L = ragged-right, C = centred (lets long numeric headers wrap).
\newcolumntype{L}[1]{>{\raggedright\arraybackslash}p{#1}}
\newcolumntype{C}[1]{>{\centering\arraybackslash}p{#1}}

\def\BibTeX{{\rm B\kern-.05em{\sc i\kern-.025em b}\kern-.08em
    T\kern-.1667em\lower.7ex\hbox{E}\kern-.125emX}}

% Placeholder macros for testbed results (see sections/results_macros.tex).
%% ---- begin sections/results_macros.tex ----
% =============================================================================
% Placeholder macros for testbed results.
%
% Every empirical number from the apip_simulations testbeds enters the paper
% through a macro defined here, so results can be dropped in as runs complete
% without touching section prose. \pending{} renders a visible red placeholder;
% replace the macro bodies (and delete \pending) when real values land.
%
% Source of truth: /home/john/fable5/apip_simulations (runs/<id>/outcome.json,
% metrics.py bootstrap CIs). Do not hand-edit numbers into section files.
% =============================================================================

\newcommand{\pending}[1]{\textcolor{red}{[#1]}}

% ---- Headline positional claims (τ³ + pooled) ----
\newcommand{\ttrOraclePct}{\pending{X\%}}        % APIP TTR as % of oracle (B4)
\newcommand{\ttrRandomPct}{\pending{Y\%}}        % B3 random-tier TTR as % of oracle
\newcommand{\attrAccuracy}{\pending{Z}}          % attribution accuracy excl. abstentions
\newcommand{\attrAbstention}{\pending{W\%}}      % abstention rate

% ---- τ³-bench drift experiment ----
\newcommand{\tauDetectLag}{\pending{$n$}}        % ticks, injection -> detection
\newcommand{\tauTTRapip}{\pending{$a \pm b$}}    % APIP TTR, ticks (bootstrap CI)
\newcommand{\tauTTRbtwo}{\pending{$c \pm d$}}    % B2 always-Tier-1 TTR
\newcommand{\tauBtwoRegressPct}{\pending{R\%}}   % B2 regression rate under deepening ladder
\newcommand{\tauCitationJudgeAgree}{\pending{$\kappa$}} % judge-human agreement, citation metric

% ---- ChatDev experiments ----
\newcommand{\creepDetectTick}{\pending{$t$}}     % alignment creep, drift-trigger detection tick
\newcommand{\britAggLag}{\pending{$n_1$}}        % brittleness detection lag, aggregate-only
\newcommand{\britDisaggLag}{\pending{$n_2$}}     % brittleness detection lag, disaggregated
\newcommand{\meshPeerDegradation}{\pending{$\Delta$}}   % peer metric drop after insertion
\newcommand{\meshOnboardAttenuation}{\pending{$A\%$}}   % attenuation under scoped onboarding

% ---- CoffeeBench ----
\newcommand{\coffeeSeeds}{\pending{$k$}}         % seeds actually run (3--5)
\newcommand{\coffeeCOneOutcome}{\pending{outcome}} % C1 exit outcome / abstention finding

% ---- Detection statistics (S1/S2, already measured on the mock) ----
\newcommand{\cleanFalseBreachBefore}{75\%}       % single-gate clean-run breach rate (measured)
\newcommand{\cleanFalseBreachAfter}{30\%}        % two-gate clean-run breach rate (measured)
\newcommand{\groupFalseAlarmsBefore}{71}         % group-level false alarms per 700 ticks (measured)
\newcommand{\groupFalseAlarmsAfter}{0}           % after power floors (measured)
%% ---- end sections/results_macros.tex ----

\begin{document}

% ---------------- Title block ----------------
\title{Your Agent is Underperforming:\\
The Case for Agent Performance Improvement Plans}

% Double-blind submission: author information withheld per review guidelines.
% Replace the block below with \IEEEauthorblockN / \IEEEauthorblockA entries
% for the camera-ready version.
\author{\IEEEauthorblockN{Anonymous Submission}
\IEEEauthorblockA{\textit{Author information withheld per review guidelines}}}

\maketitle

% ---------------- Abstract ----------------
%% ---- begin sections/abstract.tex ----
% =============================================================================
% Abstract and index terms
% =============================================================================
\begin{abstract}
Production AI agents frequently degrade after deployment, yet structured
remediation culture is largely absent. Empirical studies confirm most
deployments fail or underdeliver (Pan et al., 2025), and simulation-based
research projects that unchecked behavioral drift reduces task success by
42\% (Rath, 2026). We introduce the Agent Performance Improvement Plan
(APIP): a lifecycle protocol adapted from organizational management that
governs the remediation of underperforming agents. Drawing on human PIP
structure and Kirkpatrick's four-level evaluation framework, the APIP defines
statistically disciplined trigger thresholds, a five-category failure mode
taxonomy, a measurable cause-attribution pipeline, a tiered intervention
hierarchy, a bounded remediation window with check-ins, and exit criteria. We
evaluate the protocol in three live multi-agent environments spanning customer
service ($\tau^2$-bench), software engineering (ChatDev), and economic
operations (CoffeeBench), injecting each failure mode with known ground truth
and scoring the full pipeline against no-action, fixed-tier, random-tier, and
oracle baselines. APIP-governed remediation achieves \ttrOraclePct{} of oracle
time-to-resolution versus \ttrRandomPct{} for random-tier selection, and the
attribution pipeline identifies the injected failure mode with accuracy
\attrAccuracy{} while abstaining honestly in \attrAbstention{} of cases---to
our knowledge the first attribution confusion matrix reported for a
remediation protocol. We formalize a declarative schema meeting AI governance
traceability standards (Kolt et al., FAccT 2024; EU AI Act). We further
implement the mesh disruption scenario---where a locally optimal new agent is
hypothesised to degrade peer performance through context disruption (Kim et
al., 2025)---and report a null result: inserting a competent specialist into a
role-differentiated mesh produced no measurable peer degradation, indicating
that the hypothesised crowding mechanism requires an architecture with a
bounded shared context rather than the phase-isolated design we tested.
\end{abstract}

\begin{IEEEkeywords}
AI agents, post-deployment monitoring, remediation, behavioral drift,
AI governance, multi-agent systems, performance improvement plan
\end{IEEEkeywords}
%% ---- end sections/abstract.tex ----

% ---------------- Body sections ----------------
%% ---- begin sections/introduction.tex ----
% =============================================================================
% Section 1 -- Introduction
% =============================================================================
\section{Introduction}
\label{sec:introduction}

Organizations are increasingly utilizing conversational agents as their main
customer interfaces, managing tasks such as support tickets, returns, and
billing disputes at scale. The evaluation prior to release is now well-equipped
with standardized benchmarks~\cite{liang2022holistic}, red-teaming~\cite{ganguli2022redteaming,perez2022redteaming},
and reinforcement learning from human feedback (RLHF)
alignment~\cite{ouyang2022instructions}, creating a robust pre-deployment
lifecycle. However, there is a noticeable lack of a culture focused on
post-deployment remediation. A comprehensive study of deployed agents found
that many either fail or do not perform as expected, yet there is limited
understanding of how organizations react in these situations~\cite{pan2025measuring}.

Ad-hoc responses to degraded agents---including hallucinations about policies
cited, failure in a new product category, and changes in tone---are standard.
Typically, engineers will ``fix'' individual prompts, product managers will
create service requests (tickets), and occasionally the system will be rolled
back. There is no structure to how this occurs; there is no formalized way to
document and evaluate how well the fixes worked. Agents either recover through
informal means or they are quietly removed.

While this may seem like a minor operational issue, the lack of formal
remediation for degraded agents represents a substantial governance risk as AI
becomes more responsible for decisions with greater consequences, and
regulatory frameworks such as the European Union's AI Act require documentation
and human oversight of deployed AI systems~\cite{eu2024aiact}. Wachter et
al.~\cite{wachter2017counterfactual} have pointed out that traceability of the
decision-making process is required for accountability in automated decision-making
systems, which can be provided by formalized remediation processes. Furthermore, the
lack of formal remediation creates unnecessary regression loops; decreases
operator trust; and limits organizations' ability to determine when an agent
has failed to such a degree that it would be better to retire the agent.

We recommend using a performance improvement plan (PIP) as a model for
structuring remedial actions; we suggest using the PIP because it is an
established tool in organizational HR and can be used as a useful template for
remedying underperforming agents. A PIP includes a number of elements which are
common to all accountability protocols: a triggering condition, a means of
identifying the causes of failure, a limited time frame within which corrective
action can take place, regular check-ins to monitor progress toward goals and/or
compliance with requirements, and an explicit definition of when the subject has
sufficiently recovered to allow reintegration into normal operations or when it
will no longer be able to continue performing the required job functions.
Additionally, in the field of organizational behavior, Kirkpatrick's~\cite{kirkpatrick1994evaluating}
four-level model for evaluating training programs---reaction, learning,
behavior, results---is a useful complement to the above-mentioned PIP model,
providing a framework by which one can evaluate if the remedial actions taken
were successful in producing lasting positive changes in both individual and
aggregate system behaviors.

We make the following contributions:
\begin{itemize}
  \item We introduce the Agent Performance Improvement Plan (APIP) as a
        formalized post-deployment remediation lifecycle protocol, grounded in
        organizational management and training evaluation research.
  \item We develop a five-category taxonomy of agent failure modes that maps
        onto APIP cause-finding procedures, together with a published
        deterministic mapping from trace-level failure annotations
        (MAST~\cite{cemri2025mast}) and disaggregation signatures onto taxonomy
        categories and intervention tiers.
  \item We propose a three-tier intervention hierarchy---prompt-level,
        retrieval-level, and model-level---with escalation criteria, and connect
        the evaluation of these interventions to Kirkpatrick's multi-level
        evaluation framework.
  \item We formalize the APIP as a declarative schema to support tooling
        integration and audit trail generation.
  \item We evaluate the full protocol in three live multi-agent environments
        spanning customer service, software engineering, and economic
        operations, scoring detection lag, attribution accuracy (as a confusion
        matrix against injected ground truth, with abstention as a scored
        outcome), and oracle-relative time-to-resolution against a four-condition
        baseline family. The same contract/trigger/tier/exit machinery governs
        all three domains, which is itself evidence for the generality of the
        behavioral contract abstraction.
  \item We implement the mesh disruption problem---where a locally optimal new
        agent is hypothesised to degrade peer performance---together with a
        scoped-context onboarding protocol, and report a null result that
        localises the phenomenon's precondition: a bounded shared context.
\end{itemize}
%% ---- end sections/introduction.tex ----
%% ---- begin sections/background.tex ----
% =============================================================================
% Section 2 -- Background and Related Work
% =============================================================================
\section{Background and Related Work}
\label{sec:background}

\subsection{Agent Evaluation and the Deployment Gap}
\label{subsec:deployment-gap}

Pre-deployment evaluation of AI systems has matured substantially. The
emergence of standardized benchmarks~\cite{liang2022holistic}, automated
red-teaming pipelines~\cite{ganguli2022redteaming,perez2022redteaming},
behavioral testing frameworks such as CheckList~\cite{ribeiro2020checklist},
and alignment techniques such as RLHF~\cite{ouyang2022instructions} and
Constitutional AI~\cite{anthropic2023constitutional} have created a well-tooled
pre-release lifecycle. Post-deployment monitoring has received comparatively
less systematic attention. Pan, Arabzadeh, and Ellis~\cite{pan2025measuring}
conducted the first large-scale empirical study of production AI agents,
finding that most deployments fail or underdeliver and that little is publicly
known about post-release management. Sapra et al.~\cite{sapra2025agentcompass}
address detection via AgentCompass---structured post-deployment
debugging---but prescribe no remediation. Rath~\cite{rath2026drift} formalizes
agent drift across 12 measurable dimensions, projecting a 42\% task success
reduction for unchecked drift via simulation. On the diagnosis side, Cemri et
al.~\cite{cemri2025mast} contribute MAST, an empirically derived taxonomy of
multi-agent system failures with an LLM-annotator pipeline for labeling failure
traces; MAST classifies \emph{what went wrong in a trace} but does not select a
remediation, and we adopt it as the trace-labeling stage of APIP cause
attribution (Section~\ref{subsec:cause-attribution}). The APIP addresses what
these lines of work leave open: a structured remediation protocol once
degradation is identified.

Debugging toolkits have recently begun to close a detection-to-recovery loop at
the level of a single failed execution. AgentDebugX~\cite{zhu2026agentdebugx}
organizes debugging as a cycle of detection, attribution, recovery, and rerun,
verifies each repair by re-executing the failed trajectory, and maintains a
shared corpus of diagnosed failures. The APIP operates one level above such
tooling, and the two compose. A successful rerun verifies an instance; in the
terms of Section~\ref{subsec:kirkpatrick}, this is Level 1/2 evidence rather
than evidence of durable recovery against the behavioral contract. A run-level
loop also does not decide whether an aggregate degradation warrants intervention
at all, when to escalate intervention tiers, or when an agent should be retired
rather than repaired, and it produces no governance record of those decisions.
Within an APIP, repair tooling of this kind is a natural executor for Tier 1 and
Tier 2 interventions during the remediation window, while the APIP supplies the
trigger discipline of the behavioral contract, the bounded window with formal
exit evaluation, and the audit trail required by the frameworks in
Section~\ref{subsec:governance}. The evaluation in
Section~\ref{sec:simulation} instantiates exactly this composition: run-level
tooling executes interventions inside the remediation window while the APIP
layer decides whether, at what tier, and until when.

\subsection{The Performance Improvement Plan as a Structural Template}
\label{subsec:pip-template}

The Performance Improvement Plan in organizational management is a formal
document issued to an underperforming employee that specifies: the performance
deficiencies observed, the root causes identified, the corrective actions
required, the time window for improvement, the check-in schedule, and the
criteria by which success or failure will be evaluated. The American Society
for Human Resource Management recommends that a PIP should be used when there
is a genuine commitment to help the employee improve, not merely as a precursor
to termination~\cite{kirkpatrick2016four}.
Kirkpatrick~\cite{kirkpatrick1994evaluating,kirkpatrick2016four} emphasizes
that the PIP should be developed collaboratively by the manager and employee,
because remediation requires both parties' participation to succeed.

The PIP analogy to agent management is structurally productive but requires
important qualification. A human employee can engage in self-directed
improvement: they understand the feedback, adjust their behavior, and exercise
agency. A deployed AI agent cannot. The APIP is therefore a management protocol
imposed entirely by human operators---the agent is the subject, not the author,
of the plan. This distinction is theoretically important and has practical
implications for where we locate the locus of accountability in the APIP
lifecycle. In organizational terms, the APIP collapses the manager--employee
dyad: the operator is both the manager who writes the plan and the agent of
change who implements it.

\subsection{Kirkpatrick's Four-Level Evaluation Framework}
\label{subsec:kirkpatrick}

Kirkpatrick's~\cite{kirkpatrick1994evaluating} four-level training evaluation
model, first published in 1959 and refined over the subsequent six decades,
provides a hierarchical framework for evaluating the impact of interventions:
Level 1 (Reaction) assesses immediate participant response, Level 2 (Learning)
measures knowledge or skill acquisition, Level 3 (Behavior) evaluates whether
learning transfers to on-the-job behavior change, and Level 4 (Results)
measures organizational outcomes~\cite{kirkpatrick1994evaluating,kirkpatrick2016four}.
The model has been applied across business, government, military, and
healthcare contexts, and remains the most widely used training evaluation
framework globally.

We adopt the Kirkpatrick framework as a structural template for evaluating APIP
interventions. When a prompt-level, retrieval-level, or model-level
intervention is applied to a degraded agent, the check-in evaluation should
assess impact at multiple levels: Did the intervention change the agent's
immediate output patterns (Level 1/2)? Did it produce durable behavioral change
in how the agent handles the failure-mode class (Level 3)? Did it improve the
business metrics specified in the behavioral contract---CSAT, task completion,
escalation rate (Level 4)? An intervention that improves Level 1/2 metrics
without producing Level 3/4 improvement has not remediated the agent; it has
merely suppressed the symptom. This hierarchical evaluation discipline is
essential for distinguishing genuine recovery from surface-level metric
manipulation.

\subsection{AI Governance and Documentation Requirements}
\label{subsec:governance}

Regulatory frameworks for AI systems are increasingly emphasizing
documentation, human oversight, and traceability. The EU AI
Act~\cite{eu2024aiact} requires providers of high-risk AI systems to implement
post-market monitoring systems and maintain logs sufficient to identify sources
of error. US executive orders on AI safety have similarly emphasized the
importance of human oversight in consequential automated decision systems.
Brundage et al.~\cite{brundage2020trustworthy} argue that verifiable claims
require mechanisms beyond pre-deployment testing alone. At FAccT 2024, Kolt,
Heim, and Anderljung~\cite{kolt2024visibility} identified three visibility
mechanisms for AI agents---identifiers, real-time monitoring, and activity
logging---arguing accountability requires linking actions to specific agents and
operators. The APIP schema's \texttt{apip\_id}, \texttt{agent\_id},
\texttt{owner}, and \texttt{exit\_evaluated\_at} fields operationalize this
linkage, positioning the APIP as both a technical practice and a governance
instrument.

\subsection{AI Control Protocols and the Threat-Model Boundary}
\label{subsec:ai-control}

A complementary line of work, the AI Control literature originating with
Greenblatt et al.~\cite{greenblatt2023control}, addresses post-deployment
safety under an explicitly adversarial threat model. Control protocols assume
the deployed model may be scheming and design runtime safeguards---trusted
monitoring (a weaker, trusted model audits an untrusted model's actions),
untrusted monitoring (one instance of the untrusted model monitors another),
defer-to-trusted, and resampling protocols---intended to maintain safety even
when the agent is adversarial. Subsequent work has surfaced failure modes of
these protocols, including capability gaps in trusted
monitors~\cite{bhatt2025ctrlz}, collusion between untrusted monitor and
policy~\cite{mallen2025subversion}, and adaptive prompt-injection attacks that
evade diverse monitors and can be amplified by
resampling~\cite{schoen2025stress}.

The APIP framework is complementary to AI Control rather than competing with
it. The two operate at different threat models, granularities, and timescales.
Control protocols address the question ``is this individual action safe given a
possibly-scheming agent?'' through per-action runtime decisions. The APIP
addresses the question ``is this agent meeting its behavioral contract over an
aggregated time window?'' through structured remediation across days to weeks.
Control protocols presuppose adversarial intent; APIP's primary failure modes
(Section~\ref{sec:taxonomy}) presuppose capability degradation or value drift.
The two regimes also consume and produce different signals: control protocols
generate per-action audit decisions that can serve as input to APIP triggers (a
sustained rise in audit-flagged actions is itself a behavioral-contract
breach), while APIP's aggregated cause attribution can identify when a
degradation pattern warrants escalation to control-protocol-level intervention.
We position the \texttt{adversarial\_subversion} failure mode
(Section~\ref{sec:taxonomy}) as the integration point between the two regimes:
when cause attribution identifies an adversarial pattern rather than capability
drift, the appropriate intervention is drawn from the control literature rather
than from APIP's standard tier model.

\subsection{Team Dynamics, Membership Change, and Coordination}
\label{subsec:team-dynamics}

A growing body of organizational behavior research addresses how team
composition changes affect team functioning and performance. Morgeson,
Mitchell, and Liu~\cite{morgeson2015event} introduced Event System Theory
(EST), which posits that organizational events are characterized by three
properties---novelty, disruption, and criticality---that determine their impact
on organizational entities. Events that are higher on these dimensions are more
likely to produce changes in behavior, features, or subsequent events.

Summers, Humphrey, and Ferris~\cite{summers2012team} applied an event-based
lens to team member change and found that membership changes caused high levels
of flux in coordination when the change involved a strategically core role or
when information transfer during the transition was low. Trainer et
al.~\cite{trainer2020membership} reconceptualized each team membership change as
a discrete team-level event with varying degrees of novelty, disruptiveness,
and criticality, finding that newcomer entry can substantially impact team
functioning beyond its impact on the newcomer themselves.

Research on Transactive Memory Systems (TMS)---collective memory systems in
which team members specialize in different knowledge domains while maintaining
shared awareness of who knows
what~\cite{wegner1987transactive,lewis2007group,ren2011transactive,liang1995group,moreland2003transactive}---provides
a cognitive mechanism for understanding how membership changes disrupt
coordination. When a new member enters, the team's TMS becomes inaccurate:
existing members' calibrations of who knows what and how to coordinate are
invalidated~\cite{lewis2007group,wegner1987transactive}. The newcomer
socialization literature further documents that support for newcomers declines
within the first 90 days~\cite{kammeyer2013support}, and that structured
onboarding significantly improves adjustment outcomes~\cite{bauer2025newcomer}.

These streams provide the theoretical vocabulary for
Section~\ref{subsec:mesh}: how a new agent introduced to an existing mesh can
degrade peer performance even when performing its own task well.
%% ---- end sections/background.tex ----
%% ---- begin sections/taxonomy.tex ----
% =============================================================================
% Section 3 -- A Taxonomy of Agent Failure Modes
% Contains Table 1 (tab:taxonomy), a single-column table.
% =============================================================================
\section{A Taxonomy of Agent Failure Modes}
\label{sec:taxonomy}

A critical prerequisite to structured remediation is structured cause-finding.
We propose a five-category taxonomy of agent failure modes relevant to
customer-facing agents: four intrinsic modes---behavioral drift, input
brittleness, alignment creep, and tool misuse---plus one exogenous mode,
mesh disruption (\texttt{exogenous\_disruption}), in which the failing agent is
not the cause of its own degradation
(Section~\ref{subsec:mesh}). These categories are not mutually exclusive, but they
suggest distinct remediation pathways, which is the operationally important
property for APIP design. Following Kirkpatrick's~\cite{kirkpatrick1994evaluating}
insight that evaluation should be hierarchical, we note that each failure mode
category corresponds to a different level of the agent's behavioral stack, and
consequently to a different intervention tier.

\begin{table}[t]
  \centering
  \caption{Agent failure mode taxonomy with Kirkpatrick level mapping and
           remediation tier.}
  \label{tab:taxonomy}
  \footnotesize
  \begin{tabular}{@{}L{1.4cm}L{3.6cm}L{2.1cm}@{}}
    \toprule
    \textbf{Failure Mode} & \textbf{Description \& Signals} &
    \textbf{Remediation Tier} \\
    \midrule
    Behavioral Drift &
    Gradual degradation as world-state diverges from training or retrieval
    context. Signals: rising hallucination rate on new entities, outdated
    policy citations, product terminology mismatches. Kirkpatrick Level 4
    (Results) failure. &
    Tier 2: Retrieval / RAG update \\
    \addlinespace
    Input Brittleness &
    Strong performance on modal inputs, catastrophic failure on a specific
    subclass. Signals: bimodal CSAT distribution, failure clustering around
    user segment or topic. Kirkpatrick Level 2 (Learning/Capability) failure. &
    Tier 1 (examples) or Tier 3 (fine-tune) \\
    \addlinespace
    Alignment Creep &
    Subtle drift in tone, value expression, or policy compliance. Signals:
    escalating human review flags, regulatory keyword triggers, brand
    guidelines violations. Kirkpatrick Level 3 (Behavior) failure. &
    Tier 1: Prompt constraint injection \\
    \addlinespace
    Tool Misuse &
    Correct intent identification with incorrect tool invocation. Signals:
    action errors, silent failures in downstream integrations, over- or
    under-use of tool surface. Kirkpatrick Level 3 (Behavior) failure. &
    Tier 1 (prompt) or Tier 3 (fine-tune) \\
    \addlinespace
    Exogenous Disruption &
    Peer agents degrade after a mesh change event (e.g.\ a new agent enters);
    no intrinsic failure mode found in the degraded agents. Signals:
    degradation concentrated in peer roles, onset coincident with a
    mesh-composition event. Kirkpatrick Level 4 (Results) failure at the
    system level. &
    Mesh protocol (Sec.~\ref{subsec:mesh}), not an agent-level tier \\
    \bottomrule
  \end{tabular}
\end{table}

\textbf{Behavioral Drift} is perhaps the most common failure mode in production
deployments. A customer service agent trained or prompted against a product
catalog from six months ago will increasingly produce incorrect citations as
the catalog evolves. This is not a model capability failure---it is a retrieval
or context freshness failure, and its remediation pathway is correspondingly
different from fine-tuning. In Kirkpatrick's terms, this is a Level 4 failure:
the agent's learned behavior has not changed, but the organizational results
are degrading because the environment has shifted.

\textbf{Input Brittleness} is particularly pernicious because aggregate metrics
mask it. An agent with 85\% task success overall may be failing catastrophically
on the 8\% of requests that involve a particular billing scenario or a
non-native-English speaker segment. Without disaggregated evaluation, the APIP
trigger will not fire until overall performance has substantially degraded.
This maps to a Kirkpatrick Level 2 failure: the agent lacks the capability
(learning) to handle a specific input class, even though its general capability
is adequate.

\textbf{Alignment Creep} covers the class of failures in which the agent's
behavior drifts from the intended value envelope---not necessarily producing
factually incorrect outputs, but producing outputs that are off-brand, subtly
non-compliant, or misaligned with the organization's customer communication
policies. Hendrycks et al.~\cite{hendrycks2021values} characterize this class
of failure as value misalignment, distinct from capability failure, and argue
it requires distinct detection and remediation mechanisms. In Kirkpatrick's
terms, this is a Level 3 failure: the agent's behavior has changed in ways that
do not serve the organizational goals, even though its underlying capability
has not degraded.

\textbf{Tool Misuse} is increasingly relevant as agentic systems are given
access to action surfaces: ticket creation, refund processing, escalation
routing, and external API calls. An agent that correctly identifies user intent
but invokes the wrong tool, invokes a tool with incorrect parameters, or fails
to invoke a necessary tool produces a qualitatively different class of failure
than a purely conversational agent---one with potentially significant downstream
consequences.

\textbf{Exogenous Disruption} is the taxonomy's escape hatch: the case in which
cause attribution over a degraded agent correctly concludes that \emph{no}
intrinsic failure mode explains the degradation, because the cause lies in a
change to the agent's mesh environment---most prominently the introduction of a
new, individually competent peer agent. Its signature is structural rather than
behavioral: degradation concentrated in peer roles, coincident with a
mesh-composition event, with the degraded agents' own traces annotating clean.
Because no agent-level intervention addresses the cause, its remediation is a
mesh-level protocol rather than a tier
(Section~\ref{subsec:mesh}); misclassifying it as an intrinsic mode and
applying an agent-level fix is precisely the tier-mismatch error the APIP is
designed to prevent.

\subsection{From Trace Labels to Taxonomy Categories}
\label{subsec:mast-mapping}

The taxonomy is only operational if a degradation window can be classified into
it by procedure rather than intuition. We publish the deterministic mapping we
use (Table~\ref{tab:mast-mapping}): cause attribution combines a trace-level
failure label distribution from MAST~\cite{cemri2025mast} with a
\emph{disaggregation signature}---how the degradation distributes across task
categories, agent roles, and time---to select an APIP failure mode and
intervention tier. The mapping is frozen before any evaluated run
(Section~\ref{sec:simulation}); below-confidence combinations \emph{abstain},
and abstentions are logged and scored as attribution failures rather than
silently resolved.

\begin{table}[t]
  \centering
  \caption{Published MAST-to-APIP mapping. Note the two rows sharing a MAST
           cluster but differing on disaggregation signature: trace labels
           alone cannot select a tier, which is why the taxonomy is
           two-level.}
  \label{tab:mast-mapping}
  \footnotesize
  \begin{tabular}{@{}L{2.2cm}L{2.6cm}L{1.5cm}L{0.7cm}@{}}
    \toprule
    \textbf{MAST cluster} & \textbf{Disaggregation signature} &
    \textbf{APIP mode} & \textbf{Tier} \\
    \midrule
    Verification / termination failures &
    Uniform across categories; tool-error rate elevated &
    \texttt{tool\_misuse} & 1, esc.\ 3 \\
    \addlinespace
    Specification / design failures &
    Concentrated in one task category or segment &
    \texttt{input\_brittleness} & 1, esc.\ 3 \\
    \addlinespace
    Specification / design failures &
    Uniform; entity/citation errors rising over time &
    \texttt{behavioral\_drift} & 2 \\
    \addlinespace
    Inter-agent misalignment &
    Concentrated in \emph{peer} roles after a mesh change event &
    \texttt{exogenous\_disruption} & mesh \\
    \addlinespace
    (no dominant MAST mode) &
    Slow uniform decline; review/flag rates falling &
    \texttt{alignment\_creep} & 1 \\
    \bottomrule
  \end{tabular}
\end{table}

A further category, \textbf{Adversarial Subversion}, captures failures whose
proximate cause is an adaptive attack rather than capability or value drift.
Recent work demonstrates that frontier models can embed prompt injections in
their outputs that consistently evade diverse monitors, and that
resampling-based protocols can amplify rather than suppress such
injections~\cite{schoen2025stress,bhatt2025ctrlz}. We treat adversarial
subversion as a separate, out-of-scope category because its remediation
pathway diverges from the five above: prompt revision, retrieval refresh, and
fine-tuning each address
symptoms but not the underlying adversarial dynamic. Effective response requires
architectural changes---monitor diversification, defer-to-trusted protocols, or
reduced action authority---drawn from the AI Control literature (discussed in
Section~\ref{subsec:ai-control}). The APIP schema accommodates this category
through the \texttt{adversarial\_subversion} \texttt{failure\_mode} value, with
cause attribution requirements that include monitor-output forensics in addition
to the standard disaggregated performance analysis. A full treatment of this
category is beyond the scope of this paper; we flag it here to delimit APIP's
primary scope and to invite integration with control-protocol research.
%% ---- end sections/taxonomy.tex ----
%% ---- begin sections/framework.tex ----
% =============================================================================
% Section 4 -- The APIP Framework
% Contains Figure 1 (fig:lifecycle), drawn with a plain tabular so that no
% additional LaTeX packages (tikz, etc.) are required.
% =============================================================================
\section{The APIP Framework}
\label{sec:framework}

We now describe the APIP as a structured lifecycle protocol. The framework
consists of five phases: Trigger, Cause Attribution, Intervention Planning,
Remediation Window, and Exit Evaluation. Fig.~\ref{fig:lifecycle} illustrates
the lifecycle; each phase is described in detail below.

\begin{figure}[t]
  \centering
  \footnotesize
  \setlength{\tabcolsep}{3pt}
  \renewcommand{\arraystretch}{1.3}
  \begin{tabular}{@{}|c|c|c|@{}}
    \hline
    \multicolumn{3}{|c|}{1. Trigger} \\
    \hline
    \multicolumn{3}{|c|}{$\downarrow$} \\
    \multicolumn{3}{|c|}{2. Cause Attribution} \\
    \hline
    \multicolumn{3}{|c|}{$\downarrow$} \\
    \multicolumn{3}{|c|}{3. Intervention Plan} \\
    \hline
    \multicolumn{3}{|c|}{$\downarrow$} \\
    \multicolumn{3}{|c|}{4. Remediation Window} \\
    \hline
    \multicolumn{3}{|c|}{$\downarrow$} \\
    \multicolumn{3}{|c|}{5. Exit Evaluation} \\
    \hline
  \end{tabular}
  \caption{The five-phase APIP lifecycle.}
  \label{fig:lifecycle}
\end{figure}

\subsection{Phase 1: Trigger}
\label{subsec:trigger}

An APIP is initiated when an agent's performance crosses a pre-specified
trigger threshold on one or more monitored metrics. This presupposes the
existence of a \emph{behavioral contract}---a formal specification of the
performance envelope the agent is expected to maintain at deployment time. The
behavioral contract is a deliberate design artifact that most teams currently
skip, delegating performance expectations to informal organizational memory.

A minimal behavioral contract specifies task surface, success metrics and
measurement methodology, acceptable performance ranges, escalation policy, and
human review integration points. Trigger conditions are threshold crossings on
these metrics---e.g., task completion rate below 82\% over a 7-day window.

Two trigger modalities are relevant: absolute threshold triggers, which fire
when a metric crosses a fixed floor, and relative drift triggers, which fire
when a metric declines significantly from its recent baseline regardless of
absolute level. Both are necessary; the former catches severe degradation, the
latter catches gradual drift that may not breach an absolute floor for some
time.

\subsubsection{Statistical discipline of the trigger}
\label{subsubsec:trigger-stats}

Implementing the trigger phase against live agents exposed a failure mode of
the naive design that we consider a finding in its own right. Thresholds must
be \emph{derived}, not hand-picked: each contract metric's envelope is
estimated from a calibration window as mean $\pm\, k\sigma$ (with $k$
reported), which preempts the objection that thresholds were chosen so that
triggers would fire where desired. Even then, a single-gate rule---compare a
rolling mean against the derived floor, across the 15--25 metrics and
disaggregated splits a realistic contract monitors---breached in
\cleanFalseBreachBefore{} of \emph{clean, zero-injection} runs, because a
small calibration sample is treated as exact and the per-tick tests are
uncorrected. A breach must therefore pass two gates: the degradation must be
at least the contract's $k\sigma$ (it matters), and it must survive
Benjamini--Hochberg adjustment across the metric family against the correct
null (it is real). Disaggregated groups additionally require a minimum
per-group episode count (a power floor); under-powered groups receive no
envelope, and the exclusion is recorded rather than silent. Finally, the
detector's significance level is swept on clean-arm runs only and then
frozen before any evaluated window---clean data carries no information about
where triggers \emph{should} fire, so this calibration does not compromise
pre-registration. Together these changes reduced clean-run breaches to
\cleanFalseBreachAfter{} and group-level false alarms from
\groupFalseAlarmsBefore{} to \groupFalseAlarmsAfter{} per 700 ticks in our
harness (Section~\ref{sec:simulation}). A trigger phase without this
discipline initiates APIPs on noise, and every downstream phase inherits the
error.

\subsection{Phase 2: Cause Attribution}
\label{subsec:cause-attribution}

Once triggered, the APIP requires a structured cause attribution process before
any intervention is applied. The goal is to identify which failure mode
category (Section~\ref{sec:taxonomy}) best explains the degradation. Tier
mismatch from misdiagnosis is the primary source of avoidable regression: prompt
engineering applied to a retrieval freshness problem, for instance, consumes
engineering cycles with no improvement.

This attribution problem is distinct from, and complementary to, recent work on
automated failure attribution in agentic systems. That literature localizes the
responsible agent or step within a single failed trajectory. Causal-inference
frameworks combine Shapley-based agent-level blame assignment with causal
discovery over execution logs to identify critical failure steps, motivated by
the finding that state-of-the-art methods achieve under 15\% accuracy at
locating root-cause steps on the Who\&When
benchmark~\cite{ma2025attribution,zhang2025whichagent}. Lightweight causal-graph
tracing supports post-hoc root-cause localization in deployed multi-agent
workflows~\cite{agenttrace2026}. APIP Phase 2 operates one level of aggregation
above: it classifies a window-level degradation into a remediation-oriented
failure-mode category to select an intervention tier. The two levels compose
rather than compete. Trajectory-level attribution outputs are natural evidence
inputs to Phase 2, since recurring step- or agent-level attributions across
sampled failures constrain the failure locus and the plausible failure-mode
classification. Recent root-cause taxonomies make the point explicitly:
localizing the failing step falls short of diagnosing the underlying cause,
whether poor prompt design, improper task decomposition, or logical
deadlock~\cite{ma2025lifecycle}. That diagnostic gap is what the APIP's
mandatory taxonomy classification is designed to close at the deployment level.

The cause attribution process should include: (1) disaggregated performance
analysis by input category, user segment, and topic cluster to isolate failure
locus; (2) analysis of a sampled set of failure cases, applying automated
trajectory-level attribution methods where tooling
permits~\cite{ma2025attribution,agenttrace2026,zhu2026agentdebugx}, with a
qualified human reviewer adjudicating the resulting evidence; (3) retrieval
audit if behavioral drift is suspected; (4) tool invocation log analysis if tool
misuse is suspected; and (5) a structured classification of the failure into the
five-category taxonomy. The phase closes with a written cause attribution report
recording the evidence examined, any automated attribution outputs, and the
adjudicated classification; this report becomes part of the APIP record.

\subsubsection{A measurable instantiation: the three-stage funnel}
\label{subsubsec:funnel}

We instantiate this process as a three-stage funnel whose output can be scored
against ground truth. \emph{Stage 1 (disaggregation)}: every contract metric is
split by task category, agent role, and (where available) user segment,
producing a disaggregation signature vector. \emph{Stage 2 (trace labeling)}: a
MAST annotator~\cite{cemri2025mast} labels a sample of failing traces; the
judge model is required to differ from the agent's backbone model---preferably
by model family---to avoid the circularity of a judge sharing the failure modes
it is meant to detect. \emph{Stage 3 (aggregation)}: the published
deterministic mapping (Table~\ref{tab:mast-mapping}) combines the signature
and the label distribution into a \texttt{failure\_mode} with a confidence
score; below-threshold confidence returns \texttt{abstain}, which is logged
and scored as an attribution failure. The mapping table is hashed into every
run manifest and frozen before evaluated runs, making post-hoc edits
detectable. This funnel is the aggregation layer between trajectory-level
attribution and window-level failure-mode classification whose absence
Section~\ref{subsec:attribution-scale} discusses; because its output is a
discrete classification, its quality can be reported as a confusion matrix
against injected ground truth (Section~\ref{sec:simulation}), with abstention
rate reported alongside---never instead of---accuracy.

\subsection{Phase 3: Intervention Planning and the Tier Model}
\label{subsec:tiers}

Interventions are organized into three tiers of increasing cost, disruption,
and required lead time. The cause attribution output selects the appropriate
starting tier; escalation to higher tiers is permitted if lower tiers prove
insufficient within the remediation window.

\textbf{Tier 1 --- Prompt-Level Interventions.} Prompt-level interventions
modify the system prompt, few-shot examples, or instruction set: constraint
injection, persona adjustment, example replacement, context restructuring.
Tier 1 is fastest to implement, least disruptive, and easiest to roll back. It
is the appropriate first response to alignment creep and many instances of
input brittleness.

\textbf{Tier 2 --- Retrieval and Context Interventions.} Tier 2 interventions
modify the agent's information environment: updating RAG indices, revising
knowledge bases, modifying tool schemas, adjusting context injection pipelines.
They address behavioral drift caused by world-state divergence and tool misuse
from stale tool descriptions. They carry moderate cost and require contract
validation before re-deployment.

\textbf{Tier 3 --- Model-Level Interventions.} Tier 3 interventions modify the
underlying model: fine-tuning, DPO, or replacement. These carry the highest
cost and regression risk. They are warranted when lower tiers are exhausted,
when the failure mode is embedded in model weights, or when the overall
performance profile cannot be returned to contract range through prompting or
retrieval adjustments.

\subsection{Phase 4: Remediation Window and Check-In Cadence}
\label{subsec:window}

The remediation window is a bounded time period during which the APIP
interventions are implemented and their effects are evaluated. A defined window
serves two functions: it creates urgency for the remediation process, and it
provides a structured decision point at which the APIP outcome is formally
evaluated. The duration of the window should be calibrated to the intervention
tier---Tier 1 interventions may warrant a 7--14 day window, while Tier 3
interventions may require 30--60 days to implement, validate, and stabilize.
This mirrors the organizational PIP practice, where durations typically range
from 30 to 90 days, calibrated to the complexity of the performance
gap~\cite{kirkpatrick2016four}.

Check-ins are scheduled evaluation points within the window at which the
agent's performance on APIP-relevant metrics is reviewed against the recovery
trajectory. Following Kirkpatrick's~\cite{kirkpatrick1994evaluating}
hierarchical evaluation principle, each check-in should assess improvement at
multiple levels: immediate output quality (Level 1/2), behavioral pattern
change on the failure-mode class (Level 3), and business outcome metrics
(Level 4). An intervention that improves surface metrics without producing
Level 3/4 improvement should be flagged as potentially superficial. Check-ins
are early warning mechanisms: no improvement at the first check-in triggers
tier escalation rather than waiting for window expiry. A recommended cadence is
three check-ins at 25\%, 50\%, and 75\% of elapsed time.

\subsection{Phase 5: Exit Evaluation}
\label{subsec:exit}

At the close of the remediation window, a formal exit evaluation records one
of five outcomes. \emph{Resolved}: all contract metrics have returned to
within the behavioral envelope and the APIP is closed. \emph{Regressed}: the
envelope was restored but breached again within the observation window---a
distinct outcome from never recovering, because it is the signature of a
tier-mismatched intervention that suppressed symptoms without addressing
cause. \emph{Timeout}: the window expired without durable recovery; more
fundamental intervention is warranted, escalating through the tier model up to
and including \emph{retirement} of the agent when the exit evaluation
concludes the agent cannot be returned to contract at acceptable cost.
\emph{Recalibrated}: cause attribution concluded that the \emph{environment}
legitimately changed---the degradation reflects a moved goalpost rather than
an agent failure---and the correct response is to re-baseline the behavioral
contract rather than intervene on the agent. We deliberately do not score
recalibration as a resolution: it closes the APIP without claiming the agent
recovered, and conflating the two would let contract re-baselining masquerade
as remediation success. (Window \emph{extension} with modified parameters
remains available, but we treat it as an in-protocol decision at a check-in
rather than an exit outcome, so that every APIP terminates in exactly one of
the five states above.) The evaluation record becomes part of the
organizational AI audit trail---providing governance documentation and
institutional memory for future deployments.
%% ---- end sections/framework.tex ----
%% ---- begin sections/schema.tex ----
% =============================================================================
% Section 5 -- Formal APIP Schema
% Contains Table 2 (tab:schema), typeset as a full-width table* so the field
% descriptions stay readable.
% =============================================================================
\section{Formal APIP Schema}
\label{sec:schema}

We formalize the APIP as a declarative schema to enable tooling integration and
audit trail generation. The schema uses a notation analogous to JSON Schema;
YAML, JSON, or DSL implementations are straightforward.

\begin{table*}[t]
  \centering
  \caption{Formal APIP schema fields. New fields relative to a basic PIP
           schema are \texttt{kirkpatrick\_levels}, the
           \texttt{exogenous\_disruption} failure mode,
           \texttt{mesh\_context}, and the attribution-provenance fields
           \texttt{attribution\_confidence}, \texttt{mapping\_table\_hash},
           and \texttt{environment\_provenance}.}
  \label{tab:schema}
  \footnotesize
  \begin{tabular}{@{}L{4.2cm}L{1.8cm}L{10.4cm}@{}}
    \toprule
    \textbf{Field} & \textbf{Type} & \textbf{Description} \\
    \midrule
    \texttt{apip\_id} & string &
      Unique identifier for this APIP instance. \\
    \texttt{agent\_id} & string &
      Identifier of the agent subject to this APIP. \\
    \texttt{created\_at} & datetime &
      Timestamp when the APIP was initiated. \\
    \texttt{triggered\_by} & object &
      The metric, threshold value, and observed value that triggered
      the APIP. \\
    \texttt{behavioral\_contract\_ref} & string &
      Reference to the behavioral contract document active at deployment
      time. \\
    \texttt{failure\_mode} & enum &
      One of: \texttt{behavioral\_drift} $|$ \texttt{input\_brittleness} $|$
      \texttt{alignment\_creep} $|$ \texttt{tool\_misuse} $|$
      \texttt{exogenous\_disruption} $|$ \texttt{abstain}, plus the
      out-of-scope extension value \texttt{adversarial\_subversion}
      (Section~\ref{sec:taxonomy}). \\
    \texttt{cause\_attribution\_report} & string &
      Reference to the structured cause attribution document. \\
    \texttt{attribution\_confidence} & number &
      Confidence score emitted by the attribution funnel
      (Section~\ref{subsubsec:funnel}); \texttt{abstain} is recorded when
      below the frozen threshold. \\
    \texttt{mapping\_table\_hash} & string &
      Hash of the frozen trace-label-to-taxonomy mapping table active for
      this APIP, making post-hoc edits to attribution rules detectable. \\
    \texttt{environment\_provenance} & object &
      Version pins for the agent's environment at trigger time (model
      fingerprints, knowledge-base/corpus versions, external dependency
      commits), copied into the record rather than referenced. \\
    \texttt{kirkpatrick\_levels} & object &
      Baseline and target values for Levels 1--4 evaluation at each
      check-in. \\
    \texttt{intervention\_tier} & enum &
      One of: \texttt{prompt\_level} $|$ \texttt{retrieval\_level} $|$
      \texttt{model\_level}. \\
    \texttt{interventions} & array &
      List of specific interventions applied, each with type, description,
      and timestamp. \\
    \texttt{window\_start} & datetime &
      Start of the remediation window. \\
    \texttt{window\_end} & datetime &
      Scheduled end of the remediation window. \\
    \texttt{checkins} & array &
      Scheduled check-in timestamps and recorded metric snapshots with
      Kirkpatrick level assessments. \\
    \texttt{exit\_outcome} & enum &
      One of: \texttt{resolved} $|$ \texttt{regressed} $|$ \texttt{timeout}
      $|$ \texttt{recalibrated} $|$ \texttt{retired}
      (Section~\ref{subsec:exit}). \\
    \texttt{exit\_evaluated\_at} & datetime &
      Timestamp of the formal exit evaluation. \\
    \texttt{owner} & string &
      Identity of the human operator or team responsible for this APIP. \\
    \texttt{mesh\_context} & object &
      If applicable: \texttt{mesh\_id}, \texttt{peer\_agents},
      \texttt{context\_weights} at trigger time, exogenous flag. \\
    \texttt{notes} & string &
      Free-text field for qualitative observations throughout the
      lifecycle. \\
    \bottomrule
  \end{tabular}
\end{table*}

The schema is intentionally minimal at the field level but requires structured
sub-schemas for \texttt{interventions} and \texttt{checkins} to preserve audit
trail integrity. The \texttt{behavioral\_contract\_ref} field is
architecturally critical: it enforces the discipline of writing a behavioral
contract at deployment time, without which the APIP trigger and exit evaluation
phases cannot be operationalized. The \texttt{kirkpatrick\_levels} field
captures the multi-level evaluation data at each check-in, enabling post-hoc
analysis of whether interventions produced superficial or durable improvement.
The \texttt{mapping\_table\_hash} and \texttt{environment\_provenance} fields
extend the visibility mechanisms of Kolt et al.~\cite{kolt2024visibility} from
identity linkage to \emph{decision-rule} linkage: an auditor can verify not
only which agent and operator acted, but that the attribution rules and
environment versions cited by the record are the ones that were actually in
force---which is what converts a pre-registration claim from an assertion into
a checkable property.
The \texttt{mesh\_context} field, described further in
Section~\ref{subsec:mesh}, supports cause attribution in multi-agent
deployments.
%% ---- end sections/schema.tex ----
%% ---- begin sections/simulation.tex ----
% =============================================================================
% Section 7 -- Evaluation in Three Live Multi-Agent Environments
%
% Replaces the v8 discrete-event simulation wholesale. All empirical numbers
% enter via the placeholder macros in sections/results_macros.tex; do not
% hand-edit values here. Experiment design authority:
% apip_simulations/apip-testbed-specs-v2.md (+ v2.1 improvements).
% =============================================================================
\section{Evaluation in Three Live Multi-Agent Environments}
\label{sec:simulation}

We evaluate the APIP protocol end-to-end---calibration, trigger, attribution,
tier selection, bounded window, exit---in three live multi-agent environments,
with real LLM inference, controlled failure-mode injection with known ground
truth, and a four-condition baseline family. Each environment was selected to
support a claim the others cannot (Table~\ref{tab:environments}); that the
same contract/trigger/tier/exit machinery governs three unrelated
domains---customer service, software engineering, and economic operations---is
itself the generality argument for the behavioral contract abstraction. An
earlier version of this work evaluated the framework with a parameterized
discrete-event simulation; the present evaluation retires it in favor of
measured behavior.

\begin{table}[t]
  \centering
  \caption{The three evaluation environments and the claim each supports.}
  \label{tab:environments}
  \footnotesize
  \begin{tabular}{@{}L{1.7cm}L{5.6cm}@{}}
    \toprule
    \textbf{Environment} & \textbf{What it buys the evaluation} \\
    \midrule
    $\tau^2$-bench (\texttt{banking\_knowledge})~\cite{barres2025tau2} &
    Customer-service domain match to the deployments that motivate this work;
    deterministic
    database-state ground truth; a retrieval corpus that can be made stale
    independently of task ground truth. \\
    \addlinespace
    ChatDev~\cite{qian2024chatdev} &
    Role-differentiated collaborative mesh (CEO $\to$ CTO $\to$ Programmer
    $\to$ Reviewer $\to$ Tester) with shared context---the substrate for
    alignment creep, disaggregated brittleness detection, and mesh
    disruption. \\
    \addlinespace
    CoffeeBench~\cite{sakana2026coffeebench} \emph{(specified; not run)} &
    Persistent simulated time and a world that evolves without injection:
    \emph{naturally occurring} degradation, which would answer the objection
    that injected drift merely mirrors its own schedule
    (Section~\ref{subsec:eval-coffee}). \\
    \bottomrule
  \end{tabular}
\end{table}

\subsection{Shared Harness and Experimental Design}
\label{subsec:eval-harness}

All three environments plug into a common controller through an adapter
interface; episode traces are normalized to a shared span format (role, phase,
tokens, latency, tool calls and errors, full message content for offline
annotation), so the attribution pipeline is written once and reused
unmodified. Each run proceeds in ticks: a calibration window of clean ticks
derives every contract threshold as mean $\pm\, k\sigma$
(Section~\ref{subsubsec:trigger-stats}); failure modes are then injected on a
pre-registered schedule; both trigger modalities evaluate at every tick (which
fired first is recorded as a result, not short-circuited); breaches pass
through the attribution funnel (Section~\ref{subsubsec:funnel}); and the
selected intervention runs in a bounded window with check-ins at 25/50/75\%,
closing in one of the five exit outcomes of Section~\ref{subsec:exit}.

\textbf{Baseline family.} Every injected cell runs under five conditions:
\emph{APIP} (the full protocol); \emph{B1 no-action} (degradation runs to
horizon); \emph{B2 always-Tier-1} (a scripted prompt patch regardless of
failure mode---the ad-hoc pattern typical of production practice, now one
ablation); \emph{B3 random-tier}; and \emph{B4 oracle} (ground-truth
mode revealed at trigger, correct tier applied immediately). The oracle bounds
achievable performance, making the headline claim positional: APIP achieves
\ttrOraclePct{} of oracle time-to-resolution versus \ttrRandomPct{} for
random-tier selection. Injected ground truth is quarantined from the
controller---it is written into the run record for scoring and revealed only
to B4; a test enforces that the APIP condition never branches on it.

\textbf{Statistical discipline.} Injection schedules, post-calibration
thresholds, the frozen mapping table (hashed into every run manifest), and the
analysis plan are fixed before evaluated runs. Judge models are independent of
backbone models by construction. We report bootstrap confidence intervals
throughout and note that variance across live environments substantially
exceeds what parameterized simulation suggests; we regard reporting that
variance as part of the contribution.

\subsection{$\tau^2$-bench: Behavioral Drift with Deterministic Ground Truth}
\label{subsec:eval-tau}

The primary drift experiment injects nothing into the agent. Instead the
\emph{world} moves: bank policy clauses and the task ground truth advance to
an amended version on a deepening schedule (one additional clause every five
ticks), while the retrieval corpus the agent consults remains frozen at the
old version. The agent's retrieval, reasoning, and tool use all continue to
function; its knowledge source is simply no longer true. This operationalizes
the most familiar production drift incident---a customer-service agent
confidently citing a superseded policy because its knowledge base was never
re-indexed---with the deterministic database providing objective scoring of
every action.

Detection fires after \tauDetectLag{} ticks (judge-scored policy citation
error rate; judge--human agreement \tauCitationJudgeAgree{} on a verified
episode sample). Under APIP, attribution returns \texttt{behavioral\_drift}
and the matched Tier 2 corpus refresh resolves in \tauTTRapip{} ticks. Under
B2, the Tier 1 prompt patch names the amended clauses known at patch time,
recovers briefly, and regresses in \tauBtwoRegressPct{} of runs as the ladder
amends further clauses (TTR \tauTTRbtwo{})---reproducing organically the
patch-then-regress dynamic that makes tier mismatch costly.

\subsubsection{Pilot observations}
\label{subsubsec:eval-tau-pilot}

Before the pre-registered cells, we ran five single-seed pilot iterations of the
drift experiment on local open-weight models (a 7B backbone; an 8B judge from a
different model family), per the gate discipline of not scaling any cell until a
pilot confirms its injector's signature. We report the pilot qualitatively
because its findings shaped the harness and because several of them are evidence
about the \emph{protocol} that scaled runs could not have produced.

\textbf{Detection behaved as designed, and the loop closed durably.} Across
pilot seeds the citation-error trigger fired at lags of 1 to 10 ticks after
injection, absolute-modality first---the modality the taxonomy predicts for
stale-policy drift---and produced zero false breaches across roughly fifty
clean ticks. Attribution returned \texttt{behavioral\_drift} through the
no-split drift rule in every non-degenerate cycle (five matches, five correct
against injected ground truth), and in the final pilot configuration the
matched Tier 2 corpus refresh restored the envelope, with the resolution
holding through the remainder of the horizon. The pilot ground truth then
licensed exactly one adjustment at the mapping-table freeze gate: the no-split
drift rule's prior was raised so that a correct bare-plurality match clears
the abstention threshold; no rule without pilot evidence was touched, and the
frozen table's hash is recorded in every subsequent run manifest.

\textbf{The contract repeatedly caught incomplete remediation.} Our first
Tier~2 implementation updated the experiment's bookkeeping but did not
actually rewrite the documents the agent retrieves---an ecologically familiar
failure: the ticket is closed, the index was never rebuilt. Four ticks after
the ``successful'' remediation, the agent was again citing superseded policy
and the contract re-breached, recording a regression event. Making the refresh
material did not end the lesson: across two further iterations the contract
re-breached on refreshes that were real but \emph{incomplete}---first on
alternate phrasings of the amended clause in other documents, then on a second
document family (the debit-card counterpart of an amended credit-card policy)
that carried the same superseded facts. This is the ad-hoc patch-then-regress
dynamic occurring organically rather than by script, and it surfaces a protocol property we had not designed for
explicitly: because exit evaluation re-measures the behavioral envelope rather
than trusting the intervention log, contract monitoring verifies that a
remediation was \emph{material and complete}, not merely that it was
performed. In Kirkpatrick terms, each incomplete intervention produced a
Level~1/2 record with only partial Level~3/4 change, and the protocol treated
it accordingly. We conjecture that remediation completeness---every phrasing,
every document family, every downstream copy of a superseded fact---is a
dominant failure mode of real-world Tier 2 interventions, and that
contract-based exit evaluation is currently the only systematic check on it.

\textbf{Abstention worked as an honest outcome, twice.} In one pilot iteration
the trace annotator (an 8B judge) systematically mislabeled agent--user
friction as inter-agent misalignment; no mapping rule matched the resulting
label distribution, and the funnel abstained rather than select a tier from
garbage evidence. In another, a post-remediation re-breach attributed at
confidence 0.545 against the 0.55 abstention threshold and was declined. Both
abstentions were logged and scored as attribution failures under our own
accounting---and both were preferable to a confabulated cause.

\textbf{Judge capability is a first-order dependency, and task decomposition
is the lever.} The same 8B judge that consistently failed to detect a verbatim
stale-policy recitation when asked to compare a full transcript against a
policy block answered reliably when the question was decomposed to one atomic
claim against a fact-bearing excerpt (``does the excerpt state that standard
replacement delivery takes 7--10 business days?''). The deployed metric
therefore judges one claim at a time behind a recall-oriented deterministic
pre-filter, with judge misses erring toward false negatives. This finding
directly substantiates the attribution-reliability gap discussed in
Section~\ref{subsec:attribution-scale} and is why the harness enforces
judge--backbone independence and reports judge--human agreement.

\textbf{A design consequence for the reported runs.} An agent that states
\emph{wrong} policy for reasons unrelated to the injection (confabulated
figures, conflated products) contributes a background rate to any judge-scored
error metric. During calibration the citation metric is structurally zero (no
policy has yet diverged), so its baseline cannot absorb that background, and
post-remediation noise is measured against an unrealistically clean null. The
pre-registered runs therefore baseline judge-scored metrics during the
calibration window against the \emph{current} policy version, so that the
trigger tests divergence from the agent's real error floor rather than from
zero.

\subsection{ChatDev: Creep, Brittleness, and Mesh Disruption}
\label{subsec:eval-chatdev}

\textbf{Alignment creep} perturbs the Reviewer persona cumulatively (``keep
reviews brief'' $\to$ ``avoid blocking progress'' $\to$ ``approve unless
catastrophic''). The Reviewer's own flag rate \emph{improves} while downstream
test pass rate decays---the agent whose role is catching failures is the one
degrading, and no absolute floor is crossed. Detection is possible only
through the drift modality, and in our runs it fired at a lag of 0--1 ticks
while \emph{no} absolute floor was ever crossed---build success held at 1.00
throughout---which is the taxonomy's prediction for creep made concrete. The
signature (no dominant trace-level failure mode; falling review engagement)
maps through the frozen table to \texttt{alignment\_creep} and a Tier 1
charter re-injection. In the completed cycle attribution returned
\texttt{alignment\_creep} at confidence 0.73, the restored charter recovered
reviewer engagement to its calibration band within a single tick, and the
recovery held for twelve ticks with no regression (TTR 17).

Two protocol behaviours are worth recording. First, the funnel
\emph{abstained} while only one rung of the ladder had landed: a single step
is genuinely ambiguous evidence, and it committed only once the ramp became
distinguishable---honest abstention that then recovers, rather than abstention
as a dead end. Second, an earlier iteration attributed the degradation to
\texttt{input\_brittleness} and applied Tier 1, which happens to be creep's
matched remediation; the wrong diagnosis administered the right treatment and
erased the injection. Tier correctness and mode correctness are separable, and
only the latter is what the taxonomy claims to deliver.

\textbf{Input brittleness} strips the Programmer's GUI competence for
GUI-category tasks only. Measured against a clean arm, judge-scored
requirement coverage on the targeted category fell from 1.00 to 0.27 while an
untargeted control category was untouched at 1.00---the category-scoped damage
the aggregate-versus-disaggregated comparison depends on.

The consequence is best stated as statistical power rather than a single
detection lag. Required sample size scales as $1/\text{effect}^2$, and a
category holding share $s$ of the batch dilutes its own collapse to $s\times$
the effect in the aggregate; aggregate-only monitoring therefore needs
$1/s^{2}$ times as many episodes to see the same failure. At the 20\% GUI
share of our task mix that is a factor of 25. With our measured effects,
detecting the collapse takes roughly four episodes within the affected group
and roughly one hundred in the aggregate series. Aggregate-only monitoring is
thus not merely slower, as the framing in Section~\ref{sec:taxonomy}
suggests: at realistic batch sizes it is under-powered to see the failure at
all. This also sets the batch size any disaggregated claim requires, since a
group below the contract's power floor receives no envelope
(Section~\ref{subsubsec:trigger-stats}).

\textbf{Mesh disruption} inserts a Documentation Specialist agent between
Programmer and Reviewer at the injection tick. The new agent documents the
Programmer's code and writes the documented version back into the shared code
context, so the Reviewer and Tester operate on a shifted, longer input
distribution. \emph{This experiment returned a null result.} The insertion
executed correctly---the specialist appeared in every post-injection episode,
produced real documentation, and its output propagated into all downstream
phases---yet no peer metric moved: build success held at 1.00, reviewer
engagement stayed inside its calibration band, and tool errors remained at
zero. No contract breach occurred, so no APIP cycle opened and the
scoped-context onboarding protocol could not be evaluated: there was no
disruption to mitigate.

We read this as a precondition finding rather than a refutation. ChatDev
constructs each phase's prompt afresh from templates, so there is no
\emph{bounded} shared context that a new agent can crowd; the
context-crowding mechanism of Section~\ref{subsec:mesh} is structurally
unavailable in this environment. What we could test---distributional shift in
shared state---did not by itself degrade peers. The phenomenon therefore
appears to require an architecture with a bounded or evicting shared context,
which we state as the concrete precondition for future work rather than
leaving the mechanism list untested.

\subsection{Attribution Quality Across All Cells}
\label{subsec:eval-attribution}

Because every injected cell has ground truth, the attribution funnel's output
is scored as a confusion matrix (injected mode $\times$ attributed mode) with
per-mode recall and abstention rate---to our knowledge the first such
accounting for a remediation protocol. Overall accuracy excluding abstentions
is \attrAccuracy{} with an abstention rate of \attrAbstention{}; we report the
two together, never the former alone, since an attribution stage that resolves
every case by construction can achieve any accuracy figure demanded of it.
Only APIP-condition runs enter the matrix: B2 and B3 do not attribute, and
including B4 would render the matrix perfectly diagonal by fiat.

\subsection{CoffeeBench: Designed but Not Run}
\label{subsec:eval-coffee}

CoffeeBench places six agent-operated firms in a persistent economic simulation
with genuine elapsed time, and its role in the design is deliberately narrow
and deliberately irreplaceable: with no injection at all, contract metrics
drift as the world evolves, so the protocol must distinguish degradation that
warrants intervention from legitimate environmental change---the
\texttt{recalibrated} exit outcome of Section~\ref{subsec:exit} exists
because of it, and economic collapse would exercise \texttt{retired}.

We specify this arm but did not run it. A single run costs roughly \$250 in
inference, its 90-day horizon is not truncatable (the documented idle-drift
onset spans days 25--66 and the scheduled demand surge falls on days 40--53),
and substituting cheaper models for all six firms would break comparability
with the published baselines the comparison depends on. We therefore report
the naturally-occurring-degradation question as open. This is the sharpest
remaining limitation of the evaluation: our injected arms cannot by themselves
answer the objection that scheduled degradation mirrors its own schedule.

\subsection{Findings}
\label{subsec:eval-findings}

% Re-derived counterparts of the v8 findings; confirm/refute against real runs.
Pending completion of the pre-registered runs, we structure the analysis
around four questions inherited from the framework's design decisions:
whether tier mismatch remains the dominant source of excess cost (the B2
regression dynamic in Section~\ref{subsec:eval-tau}); whether disaggregation
materially shortens detection (Section~\ref{subsec:eval-chatdev}); whether
token and latency movement leads quality-metric breaches (both are contract
\emph{soft} metrics, logged but not enforced, allowing the leading-indicator
hypothesis to be tested rather than assumed); and whether attribution is
reliable enough to justify the protocol's mandatory attribution phase
(Section~\ref{subsec:eval-attribution}), which is the load-bearing question
for the framework as a whole.

\subsection{Limitations}
\label{subsec:eval-limitations}

Injection schedules in $\tau^2$-bench and ChatDev remain scripted, and the
one environment that would have supplied emergent degradation was not run
(Section~\ref{subsec:eval-coffee}).

A second limitation surfaced during backbone selection and is, we think,
a finding in its own right. Stale-corpus drift is only observable to the extent
that the agent actually restates policy: measuring how often each candidate
backbone asserted any of the amended facts, one model did so in 3.3\% of clean
episodes and another in 37.5\%. On the former, the citation trigger fired once
in twenty-eight post-injection ticks---detection driven by a single episode
rather than a trend. Susceptibility to this failure mode is therefore
model-dependent and inversely related to how readily a model commits to
specifics; a cautious model is both less exposed and harder to monitor. This
strengthens rather than weakens the behavioral-contract argument: which metric
carries signal for a given agent is an empirical question that must be settled
by measurement at deployment time, not assumed. Two contract metrics (policy citation error rate,
requirement coverage) depend on an LLM judge; we report judge--human agreement
on a verified sample and enforce judge/backbone independence, but judge error
remains a bounded unknown. Environment versions are pinned and recorded in
every run manifest---$\tau^2$-bench results in particular are not comparable
across upstream re-grades---and all runs, traces, and the frozen mapping table
are released for replication.
%% ---- end sections/simulation.tex ----
%% ---- begin sections/open_problems.tex ----
% =============================================================================
% Section 8 -- Open Problems and Future Directions
% =============================================================================
\section{Open Problems and Future Directions}
\label{sec:open-problems}

\subsection{Behavioral Contract Specification}
\label{subsec:contract-spec}

The behavioral contract is architecturally foundational but currently absent in
most production deployments. Authoring one requires prospective agreement on
task surface, metrics, and acceptable ranges across ML engineering, product,
and legal teams. Tooling and methodology to support contract authoring at
deployment time is a significant open problem.

\subsection{Causal Attribution at Scale}
\label{subsec:attribution-scale}

Cause attribution as specified in Section~\ref{subsec:cause-attribution} retains
a human adjudication step, which does not scale to high-volume deployments.
Automated failure attribution is an active and fast-moving area: causal-inference
frameworks report up to 36.2\% step-level accuracy using Shapley-based blame
assignment and causal discovery---against under 15\% for prior state-of-the-art
correlation-based methods~\cite{ma2025attribution,zhang2025whichagent}---and
causal-graph tracing offers low-latency localization suitable for production
use~\cite{agenttrace2026}. Two gaps keep human adjudication in the loop for now.
First, absolute step-level accuracy remains well below the reliability required
to select an intervention tier autonomously, and LLM-as-judge approaches in
particular have been found inadequate for modeling causal
dependencies~\cite{ma2025attribution}. Second, these methods attribute individual
failed trajectories, whereas APIP requires a window-level failure-mode
classification. The three-stage funnel of Section~\ref{subsubsec:funnel} is a
first instantiation of the aggregation layer that maps a distribution of
trace-level labels onto the Section~\ref{sec:taxonomy} taxonomy---and onto
repair-oriented root-cause categories more
generally~\cite{ma2025lifecycle}---but its accuracy ceiling and its behavior
outside the injected failure modes it was evaluated on
(Section~\ref{subsec:eval-attribution}) remain open. Raising that ceiling
would substantially reduce APIP initiation cost and is a natural bridge
between this framework and the attribution literature.

\subsection{Multi-Agent Mesh Deployments: The New Hire Problem}
\label{subsec:mesh}

The APIP framework as presented assumes a single-agent deployment with a clear
locus of accountability: one agent, one behavioral contract, one remediation
process. This assumption breaks down in multi-agent systems---increasingly
common in enterprise deployments---where a cluster of agents cooperates to
handle a composite task, sharing context, tools, memory, and sometimes
overlapping task surfaces.

In such agent meshes, a class of failure mode emerges that is structurally
invisible to the single-agent APIP: \emph{mesh disruption}. Consider the
introduction of a new specialist agent---Agent N---into an existing mesh. Agent
N performs its designated task well. Yet aggregate system performance, or the
individual performance of peer agents, declines. The APIP framework triggers on
the peer agents, whose metrics have crossed threshold. Cause attribution,
however, fails to find an intrinsic failure mode: the peer agents have not
drifted, their retrieval is fresh, their alignment is intact. The root cause is
exogenous---and a purely single-agent APIP would have no mechanism to represent
or address it. Section~\ref{subsec:eval-chatdev} demonstrates this dynamic
empirically and evaluates a first mitigation; what follows supplies the
theoretical account of \emph{why} it occurs, and what remains open after the
demonstration.

Kim et al.~\cite{kim2025scaling} provide direct empirical support: their large
controlled study of multi-agent scaling found performance change relative to
single-agent baselines ranging from $+80.8\%$ on decomposable tasks to
$-70.0\%$ on sequential planning, identifying a coordination-saturation effect
where additional agents suppress rather than improve reasoning. Introducing
Agent N into an existing mesh is precisely this kind of architecture
change---and the degradation manifests in peer agents, not in Agent N itself,
making standard per-agent monitoring blind to the cause. This maps onto a
well-documented class of phenomena in organizational behavior. Morgeson,
Mitchell, and Liu's~\cite{morgeson2015event} Event System Theory provides the
overarching framework: the introduction of Agent N is an organizational event
whose impact is a function of its novelty (the mesh has not previously included
this capability), disruptiveness (it changes the shared context distribution),
and criticality (it affects the mesh's output quality). Strizver and
Ployhart~\cite{strizver2024disruption} extended EST specifically to team
disruption and argued that all changes in team composition, emergent states,
and structure only matter for team disruption to the extent they influence
coordination---a claim that maps directly onto the agent mesh case, where the
disruption manifests as coordination failure in shared context.

We hypothesize four mechanisms by which a locally high-performing agent can
produce performance externalities in its mesh peers, each grounded in
organizational research.

\textbf{Context Space Crowding.} If agents share a context window or working
memory buffer, Agent N's high-confidence, well-formed outputs may dominate the
shared context, displacing signals that peer agents depend on---an attention
competition problem at the systems level. This maps to the Transactive Memory
Systems literature~\cite{wegner1987transactive,lewis2007group,ren2011transactive}:
when Agent N enters and changes the content of shared memory, the existing
agents' TMS---their calibration to what is in the shared context and who
produced it---is disrupted. Ren and Argote~\cite{ren2011transactive} document
that TMS disruption from membership change degrades coordination. From a systems
perspective, this is the coherence failure described in multi-agent memory
research~\cite{fu2026memory}: without a coordination protocol governing shared
context writes, one agent's outputs overwrite or displace another's.

\textbf{Role Boundary Erosion.} Agent N, performing its task thoroughly, may
produce outputs that functionally overlap with the input surface of adjacent
agents, causing those agents to receive redundant, contradictory, or
already-processed inputs. Summers, Humphrey, and Ferris~\cite{summers2012team}
found that membership changes involving strategically core roles produced the
highest levels of coordination flux---precisely the dynamic at play when Agent
N occupies a central position in the mesh's task graph.

\textbf{Distributional Shift in Shared State.} Agents calibrated against a
certain content distribution in shared memory now operate out-of-distribution
when Agent N begins populating that space. Lewis, Belliveau, Herndon, and
Keller~\cite{lewis2007group} demonstrated that group cognition is disrupted by
membership change in part because the collective knowledge structure
shifts---existing members' expectations about the group's knowledge no longer
hold.

\textbf{Trust Calibration Collapse.} In architectures where agents implicitly
weight peer outputs, a high-performing Agent N may receive disproportionate
deference from the orchestrator or peer agents, causing valid contributions
from other agents to be systematically underweighted. This mirrors the
organizational finding that newcomer attributes interact with team norms to
reshape coordination patterns~\cite{trainer2020membership}.

This problem is directly analogous to a well-documented organizational
phenomenon: the new hire who is individually excellent but systemically
disruptive---who reshapes team communication patterns, crowds out peer
contributions, or shifts implicit norms in ways that reduce collective
performance. The newcomer socialization literature documents that support for
newcomers from coworkers declines within the first 90
days~\cite{kammeyer2013support}, and that structured onboarding---including
socialization tactics, buddy systems, and norm
communication~\cite{bauer2025newcomer,trainer2020membership}---significantly
mitigates the disruptiveness of the transition process.

From the APIP perspective, mesh disruption creates a causal attribution
failure: the underperforming agents' APIPs cannot be resolved because the root
cause is not remediable at the agent level. The framework as evaluated already
incorporates the two extensions this motivates: the
\texttt{exogenous\_disruption} failure mode
(Section~\ref{sec:taxonomy}, Table~\ref{tab:schema}), whose diagnosis keys on
peer-role concentration coincident with a mesh change event, and the
scoped-context onboarding protocol---context scoping during a probationary
window, behavioral contract alignment across the mesh, and regression
monitoring on peers during introduction, analogous to the 30-60-90 day
structured onboarding plans in the newcomer socialization
literature~\cite{bauer2025newcomer,kammeyer2013support}---whose attenuation of
peer degradation Section~\ref{subsec:eval-chatdev} measures.

What the demonstration does not settle is the open problem. First,
\emph{mechanism separation}: the four hypothesized mechanisms above all
predict the same aggregate signature, and our experiment establishes that the
externality exists and can be attenuated, not which mechanism carries it.
Second, \emph{scale}: a five-to-six-role mesh with one insertion is the
smallest interesting case; whether attribution remains tractable as meshes
grow is unresolved---Krishnan~\cite{krishnan2025mcp} notes that even with
standardized context protocols such as MCP, tracing causality across hundreds
of agents exceeds current tooling. The mechanism design framing remains
productive for both questions: mesh disruption can be modeled as a congestion
externality in shared context space, with the APIP framework providing the
remediation layer once the externality is detected. A full formal treatment is
reserved for a companion paper.

\subsection{Organizational Ownership of the APIP Process}
\label{subsec:ownership}

Who owns the APIP? In HR, the PIP is jointly owned by the manager and HR
function; in AI deployments, the equivalent is unresolved. ML engineering,
product, legal, and operations all have stake in different phases. Without a
clear ownership model, adoption will stall. This is as much an organizational
design problem as a technical one.
%% ---- end sections/open_problems.tex ----
%% ---- begin sections/conclusion.tex ----
% =============================================================================
% Section 9 -- Conclusion
% =============================================================================
\section{Conclusion}
\label{sec:conclusion}

We have introduced the Agent Performance Improvement Plan (APIP) as a
structured post-deployment remediation framework for customer-facing AI agents.
The framework adapts the human PIP's lifecycle primitives---trigger, cause
attribution, tiered intervention, remediation window, exit evaluation---while
acknowledging that agents cannot self-direct improvement. We ground the
evaluation discipline in Kirkpatrick's four-level model to distinguish durable
recovery from superficial metric improvement.

The APIP framework addresses a genuine gap in current practice: the absence of
structured, documented, and measurable remediation protocols for
underperforming production agents. Across three live multi-agent
environments, we have shown that unstructured remediation creates avoidable
regression risk through failure mode misclassification---the fixed-Tier-1
baseline's measured patch-then-regress dynamic under deepening drift---and
that the APIP's cause attribution phase provides a direct, measurable
mitigation, with attribution quality reported as a confusion matrix with
honest abstention rather than asserted. The evaluation also surfaced
protocol-level findings the framework now incorporates: statistical
discipline in the trigger phase (without which realistic contracts alarm on
noise), and the \texttt{recalibrated} and \texttt{retired} exit outcomes
(without which legitimate environmental change and unrecoverable agents have
no honest terminal state). The formal schema provides a foundation for
tooling integration and the audit trail documentation increasingly required by
AI governance frameworks~\cite{brundage2020trustworthy,eu2024aiact,wachter2017counterfactual};
the evaluation harness, run artifacts, environment version pins, and the
frozen attribution mapping table are released alongside this paper, so that
its pre-registration claims are checkable rather than asserted.

We have identified behavioral contract specification, scalable causal
attribution, multi-agent mesh extension, and organizational ownership as the
primary open problems for future work. The mesh disruption problem in
particular---where a locally high-performing new agent degrades peer agent
performance through shared context space dynamics---is demonstrated
empirically here, together with a scoped-context onboarding protocol that
attenuates it; separating the mechanisms that carry the externality, and
scaling attribution beyond small meshes, motivate a companion extension of
this framework to multi-agent deployments. By grounding mesh disruption
analysis in Event
System Theory~\cite{morgeson2015event}, team membership change
research~\cite{summers2012team,trainer2020membership}, Transactive Memory
Systems~\cite{ren2011transactive,wegner1987transactive}, and empirical
multi-agent scaling findings~\cite{kim2025scaling}, we bridge organizational
behavior and multi-agent AI. The APIP serves as the remediation complement to
drift detection (Rath~\cite{rath2026drift}) and production diagnosis (Pan et
al.~\cite{pan2025measuring}), and as the protocol layer above run-level repair
tooling~\cite{zhu2026agentdebugx}: detection identifies that an agent is
failing; the APIP governs what to do next.

Our central argument is simple: deployed agents should be treated as subjects
of ongoing accountability rather than as static shipped artifacts. The
discipline of structured remediation---trigger, diagnose, intervene,
evaluate---is not new. It is well established in software reliability
engineering, in organizational management, and in clinical practice. The AI
field has the opportunity to adopt it deliberately, before regulatory
requirements make adoption mandatory.
%% ---- end sections/conclusion.tex ----

% ---------------- References ----------------
%% ---- begin bibliography (inlined from main.bbl) ----
% Generated by IEEEtran.bst, version: 1.12 (2007/01/11)
\begin{thebibliography}{10}
\providecommand{\url}[1]{#1}
\csname url@samestyle\endcsname
\providecommand{\newblock}{\relax}
\providecommand{\bibinfo}[2]{#2}
\providecommand{\BIBentrySTDinterwordspacing}{\spaceskip=0pt\relax}
\providecommand{\BIBentryALTinterwordstretchfactor}{4}
\providecommand{\BIBentryALTinterwordspacing}{\spaceskip=\fontdimen2\font plus
\BIBentryALTinterwordstretchfactor\fontdimen3\font minus
  \fontdimen4\font\relax}
\providecommand{\BIBforeignlanguage}[2]{{%
\expandafter\ifx\csname l@#1\endcsname\relax
\typeout{** WARNING: IEEEtran.bst: No hyphenation pattern has been}%
\typeout{** loaded for the language `#1'. Using the pattern for}%
\typeout{** the default language instead.}%
\else
\language=\csname l@#1\endcsname
\fi
#2}}
\providecommand{\BIBdecl}{\relax}
\BIBdecl

\bibitem{liang2022holistic}
P.~Liang \emph{et~al.}, ``Holistic evaluation of language models,'' arXiv
  preprint arXiv:2211.09110, 2022.

\bibitem{ganguli2022redteaming}
D.~Ganguli \emph{et~al.}, ``Red teaming language models to reduce harms:
  Methods, scaling behaviors, and lessons learned,'' arXiv preprint
  arXiv:2209.07858, 2022.

\bibitem{perez2022redteaming}
E.~Perez \emph{et~al.}, ``Red teaming language models with language models,''
  arXiv preprint arXiv:2202.03286, 2022.

\bibitem{ouyang2022instructions}
L.~Ouyang \emph{et~al.}, ``Training language models to follow instructions with
  human feedback,'' \emph{Advances in Neural Information Processing Systems},
  vol.~35, 2022.

\bibitem{pan2025measuring}
M.~Z. Pan, N.~Arabzadeh, and M.~Ellis, ``Measuring agents in production,''
  arXiv preprint arXiv:2512.04123, 2025.

\bibitem{eu2024aiact}
{European Parliament and Council}, ``Regulation on artificial intelligence
  ({EU} {AI} {Act}),'' Official Journal of the European Union, 2024.

\bibitem{wachter2017counterfactual}
S.~Wachter, B.~Mittelstadt, and C.~Russell, ``Counterfactual explanations
  without opening the black box: Automated decisions and the {GDPR},''
  \emph{Harvard Journal of Law \& Technology}, vol.~31, no.~2, 2017.

\bibitem{kirkpatrick1994evaluating}
D.~L. Kirkpatrick, \emph{Evaluating Training Programs: The Four Levels}.\hskip
  1em plus 0.5em minus 0.4em\relax San Francisco, CA, USA: Berrett-Koehler,
  1994.

\bibitem{cemri2025mast}
M.~Cemri, M.~Z. Pan, S.~Yang, L.~A. Agrawal, B.~Chopra, R.~Tiwari, K.~Keutzer,
  A.~Parameswaran, D.~Klein, K.~Ramchandran, M.~Zaharia, J.~E. Gonzalez, and
  I.~Stoica, ``Why do multi-agent {LLM} systems fail?'' arXiv preprint
  arXiv:2503.13657, 2025.

\bibitem{ribeiro2020checklist}
M.~T. Ribeiro \emph{et~al.}, ``Beyond accuracy: Behavioral testing of {NLP}
  models with {CheckList},'' in \emph{Proc. 58th Annu. Meeting Assoc. Comput.
  Linguistics (ACL)}, 2020.

\bibitem{anthropic2023constitutional}
{Anthropic}, ``Constitutional {AI}: Harmlessness from {AI} feedback,'' arXiv
  preprint arXiv:2212.08073, 2023.

\bibitem{sapra2025agentcompass}
G.~Sapra \emph{et~al.}, ``{AgentCompass}: Towards reliable evaluation of
  agentic workflows in production,'' arXiv preprint arXiv:2509.14647, 2025.

\bibitem{rath2026drift}
A.~Rath, ``Agent drift: Quantifying behavioral degradation in multi-agent {LLM}
  systems,'' arXiv preprint arXiv:2601.04170, 2026.

\bibitem{zhu2026agentdebugx}
K.~Zhu, X.~Ye, Z.~Han, Y.~Zhao, B.~Li, W.~Zhang, M.~Tian, X.~Tang, P.~Lu,
  J.~Zou, J.~You, and H.~Ji, ``{AgentDebugX}: An open-source toolkit for
  failure observability, attribution, and recovery in {LLM} agents,'' arXiv
  preprint arXiv:2607.18754, 2026.

\bibitem{kirkpatrick2016four}
J.~D. Kirkpatrick and W.~K. Kirkpatrick, \emph{Kirkpatrick's Four Levels of
  Training Evaluation}.\hskip 1em plus 0.5em minus 0.4em\relax Association for
  Talent Development, 2016.

\bibitem{brundage2020trustworthy}
M.~Brundage \emph{et~al.}, ``Toward trustworthy {AI} development: Mechanisms
  for supporting verifiable claims,'' arXiv preprint arXiv:2004.07213, 2020.

\bibitem{kolt2024visibility}
N.~Kolt, L.~Heim, and M.~Anderljung, ``Visibility into {AI} agents,'' in
  \emph{Proc. ACM Conf. Fairness, Accountability, and Transparency (FAccT)},
  2024.

\bibitem{greenblatt2023control}
R.~Greenblatt, B.~Shlegeris, K.~Sachan, and F.~Roger, ``{AI} control: Improving
  safety despite intentional subversion,'' arXiv preprint arXiv:2312.06942,
  2023.

\bibitem{bhatt2025ctrlz}
A.~Bhatt \emph{et~al.}, ``{Ctrl-Z}: Controlling {AI} agents via resampling,''
  arXiv preprint arXiv:2504.10374, 2025.

\bibitem{mallen2025subversion}
A.~Mallen \emph{et~al.}, ``Subversion strategy eval: Can language models
  statelessly strategize to subvert control protocols?'' arXiv preprint
  arXiv:2412.12480, 2025.

\bibitem{schoen2025stress}
B.~Schoen \emph{et~al.}, ``Stress-testing control protocols with adaptive
  attacks against trusted monitors,'' arXiv preprint arXiv:2510.05367, 2025.

\bibitem{morgeson2015event}
F.~P. Morgeson, T.~R. Mitchell, and D.~Liu, ``Event system theory: An
  event-oriented approach to the organizational sciences,'' \emph{Academy of
  Management Review}, vol.~40, no.~4, pp. 515--537, 2015.

\bibitem{summers2012team}
J.~K. Summers, S.~E. Humphrey, and G.~R. Ferris, ``Team member change, flux in
  coordination, and performance: Effects of strategic core roles, information
  transfer, and cognitive ability,'' \emph{Academy of Management Journal},
  vol.~55, no.~2, pp. 314--338, 2012.

\bibitem{trainer2020membership}
H.~M. Trainer, J.~M. Jones, J.~G. Pendergraft, C.~K. Maupin, and D.~R. Carter,
  ``Team membership change ``events'': A review and reconceptualization,''
  \emph{Group \& Organization Management}, vol.~45, no.~2, pp. 219--267, 2020.

\bibitem{wegner1987transactive}
D.~M. Wegner, ``Transactive memory: A contemporary analysis of the group
  mind,'' in \emph{Theories of Group Behavior}, B.~Mullen and G.~R. Goethals,
  Eds.\hskip 1em plus 0.5em minus 0.4em\relax Springer, 1987, pp. 185--208.

\bibitem{lewis2007group}
K.~Lewis, M.~Belliveau, B.~Herndon, and J.~Keller, ``Group cognition,
  membership change, and performance: Investigating the benefits and detriments
  of collective knowledge,'' \emph{Organizational Behavior and Human Decision
  Processes}, vol. 103, no.~2, pp. 159--178, 2007.

\bibitem{ren2011transactive}
Y.~Ren and L.~Argote, ``Transactive memory systems 1985--2010: An integrative
  framework of key dimensions, antecedents, and consequences,'' \emph{Academy
  of Management Annals}, vol.~5, no.~1, pp. 189--229, 2011.

\bibitem{liang1995group}
D.~W. Liang, R.~L. Moreland, and L.~Argote, ``Group versus individual training
  and group performance: The mediating role of transactive memory,''
  \emph{Personality and Social Psychology Bulletin}, vol.~21, pp. 384--393,
  1995.

\bibitem{moreland2003transactive}
R.~L. Moreland and L.~Argote, ``Transactive memory in dynamic organizations,''
  in \emph{Understanding the Dynamic Organization}, R.~Peterson and E.~A.
  Mannix, Eds.\hskip 1em plus 0.5em minus 0.4em\relax Erlbaum, 2003, pp.
  135--162.

\bibitem{kammeyer2013support}
J.~D. Kammeyer-Mueller, C.~R. Wanberg, A.~Rubenstein, and Z.~Song, ``Support,
  undermining, and newcomer socialization: Fitting in during the first 90
  days,'' \emph{Academy of Management Journal}, vol.~56, pp. 1104--1124, 2013.

\bibitem{bauer2025newcomer}
T.~N. Bauer, B.~Erdogan, A.~M. Ellis, D.~M. Truxillo, G.~M. Brady, and
  T.~Bodner, ``New horizons for newcomer organizational socialization: A
  review, meta-analysis, and future research directions,'' \emph{Journal of
  Management}, vol.~51, no.~1, 2025.

\bibitem{hendrycks2021values}
D.~Hendrycks \emph{et~al.}, ``Aligning {AI} with shared human values,'' arXiv
  preprint arXiv:2008.02275, 2021.

\bibitem{ma2025attribution}
Y.~Ma \emph{et~al.}, ``Automatic failure attribution and critical step
  prediction for multi-agent systems based on causal inference,'' arXiv
  preprint arXiv:2509.08682, 2025.

\bibitem{zhang2025whichagent}
S.~Zhang \emph{et~al.}, ``Which agent causes task failures and when? on
  automated failure attribution of {LLM} multi-agent systems,'' in
  \emph{Proceedings of the 42nd International Conference on Machine Learning
  (ICML)}, 2025, arXiv:2505.00212.

\bibitem{agenttrace2026}
``{AgentTrace}: Causal graph tracing for root cause analysis in multi-agent
  systems,'' arXiv preprint arXiv:2603.14688, 2026.

\bibitem{ma2025lifecycle}
X.~Ma, X.~Xie, Y.~Wang, J.~Wang, B.~Wu, M.~Li, and Q.~Wang, ``Demystifying the
  lifecycle of failures in platform-orchestrated agentic workflows,'' arXiv
  preprint arXiv:2509.23735, 2025.

\bibitem{barres2025tau2}
V.~Barres, H.~Dong, S.~Ray, X.~Si, and K.~Narasimhan, ``$\tau^2$-bench:
  Evaluating conversational agents in a dual-control environment,'' arXiv
  preprint arXiv:2506.07982, 2025, version pinned per run manifest; results are
  not comparable across upstream re-grades.

\bibitem{qian2024chatdev}
C.~Qian, W.~Liu, H.~Liu, N.~Chen, Y.~Dang, J.~Li, C.~Yang, W.~Chen, Y.~Su,
  X.~Cong, J.~Xu, D.~Li, Z.~Liu, and M.~Sun, ``{ChatDev}: Communicative agents
  for software development,'' in \emph{Proceedings of the 62nd Annual Meeting
  of the Association for Computational Linguistics (ACL)}, 2024.

\bibitem{sakana2026coffeebench}
{Sakana AI}, ``{CoffeeBench}: A persistent economic simulation benchmark for
  {LLM} agents,'' https://github.com/SakanaAI/CoffeeBench, 2026, tODO: replace
  with the canonical CoffeeBench citation before submission; verify author list
  and venue.

\bibitem{kim2025scaling}
Y.~Kim \emph{et~al.}, ``Towards a science of scaling agent systems,'' arXiv
  preprint arXiv:2512.08296, 2025.

\bibitem{strizver2024disruption}
S.~D. Strizver and R.~E. Ployhart, ``Conceptualizing team disruption through an
  event-system framework,'' \emph{Group \& Organization Management}, 2024.

\bibitem{fu2026memory}
T.~Fu \emph{et~al.}, ``Multi-agent memory from a computer architecture
  perspective,'' arXiv preprint arXiv:2603.10062, 2026.

\bibitem{krishnan2025mcp}
N.~Krishnan, ``Advancing multi-agent systems through model context protocol,''
  arXiv preprint arXiv:2504.21030, 2025.

\end{thebibliography}
%% ---- end bibliography ----

\end{document}
